\documentclass[12pt,a4paper]{article}

\usepackage[T1]{fontenc}
\usepackage{amsmath,amsthm,amssymb,mathtools}
\usepackage{geometry}
\usepackage{microtype}
\usepackage{setspace}
\usepackage{titlesec}
\usepackage{hyperref}
\usepackage{enumitem}
\usepackage{booktabs}

\geometry{a4paper,top=2.8cm,bottom=2.8cm,left=3.2cm,right=3.2cm}
\onehalfspacing

\hypersetup{colorlinks=true,linkcolor=black,citecolor=black,urlcolor=black}

% Theorem environments
\newtheorem{theorem}{Theorem}[section]
\newtheorem{lemma}[theorem]{Lemma}
\newtheorem{corollary}[theorem]{Corollary}
\newtheorem{proposition}[theorem]{Proposition}
\newtheorem{definition}[theorem]{Definition}
\theoremstyle{remark}
\newtheorem{remark}{Remark}[section]
\newtheorem{example}{Example}[section]
\theoremstyle{definition}
\newtheorem{principle}{Principle}[section]

\titleformat{\section}{\large\bfseries}{\thesection.}{0.6em}{}
\titleformat{\subsection}{\normalsize\bfseries}{\thesubsection.}{0.5em}{}

\title{\textbf{Distinction Geometry:\\
Projection, Reconstruction Defect,\\
and the Interface-Native Metric}\\[0.5em]
\large Mathematical Core}

\author{Flyxion\\[0.4em]
\small Independent Researcher}

\date{\today}

\begin{document}

\maketitle

\begin{abstract}
We develop the mathematical core of a framework in which the primitive
object of inquiry is not a hidden state space but the geometry of
distinguishability between observable projections. Three interlocking
structures are established. First, the \emph{reconstruction defect}
$\Delta_R(h) = d(h, R(\Pi(h)))$ is introduced as a unified measure of
irreversibility: the persistent failure of any map from observations back
to the generating history. We prove that irreversibility, in this sense,
is equivalent to the non-injectivity of the projection $\Pi$, and that
the minimum achievable reconstruction defect is controlled by the
information lost under projection. Second, the \emph{distinction
functional} $\mathcal{D}(P,Q) = \operatorname{JSD}(P \| Q)$ is developed
as the canonical measure of the information destroyed when two observable
interfaces are merged. Its key properties are established: interface
nativity, the entropy decomposition, the gauge-orbit characterisation,
and the operationally precise meaning of the two entropy terms. Third,
the \emph{space of observable interfaces} $(\mathcal{M}, d_{\mathrm{int}})$
is shown to be a metric space that is locally Riemannian with Fisher
information metric. Its curvature is interpreted: positive curvature
corresponds to regions where many internal models produce nearly identical
observations (overfitting, paradigm stability, institutional opacity);
negative curvature to regions of rapid observational divergence (high
falsifiability). These three structures --- reconstruction defect,
distinction functional, and interface geometry --- form a mathematical
core from which physical, cognitive, and institutional applications
follow as corollaries.
\end{abstract}

\tableofcontents
\newpage

%----------------------------------------------------------------------
\section*{Orientation}
\addcontentsline{toc}{section}{Orientation}
%----------------------------------------------------------------------

Every observation is a projection. Whatever generates an observation ---
a physical system, a cognitive process, a social institution, a measuring
device --- projects its internal configuration onto a reduced description
that can be received, recorded, and compared. The projection is in general
many-to-one. Many internal configurations produce the same observation.
What does not vary across a fiber is, by definition, invisible.

This paper takes that structure seriously as a mathematical object. It
does not ask what the hidden internal space is, or how it might be
reconstructed. It asks what geometry is imposed on the space of
observations by the existence of fibers, and what the non-invertibility
of projection implies for irreversibility, memory, and the comparison of
theories.

The answer, developed in three sections, is that the geometry of
observation is an information geometry: the natural metric on observable
interfaces is the Jensen--Shannon distance, the curvature of the resulting
space measures falsifiability and opacity, and the minimum information
destroyed by any observation is the reconstruction defect of the generating
projection. These are not analogies between different domains. They are
the same mathematical structure appearing in different applications.

The present document is the mathematical core of a larger monograph,
\emph{Distinction Geometry}. It is intended to be read first: the
applications --- to quantum gravity, thermodynamics, cognition, and
institutional theory --- are corollaries of what is proved here.


%======================================================================
\section{The Projection Setting}
%======================================================================

\subsection{Histories, Projections, and Fibers}

Throughout this paper, $M$ denotes a \emph{space of histories}: a set
whose elements are complete trajectories of a system over a time interval,
not instantaneous configurations. This choice is deliberate. The phenomena
we wish to unify --- irreversibility, memory, learning, paradigm change ---
are inherently historical. They concern how the past constrains the future
and what of the past survives into observation. A state-based formalism
would need to carry additional temporal structure explicitly; the history
formalism carries it automatically.

\begin{definition}[Projection Setting]
A \emph{projection setting} is a triple $(M, O, \Pi)$ where $M$ is a
measurable space of histories, $O$ is a measurable space of observable
outcomes, and $\Pi : M \to O$ is a surjective measurable map called the
\emph{projection}. The \emph{fiber} over $o \in O$ is
\begin{equation}
  F_o = \Pi^{-1}(o) = \{h \in M : \Pi(h) = o\}.
\end{equation}
Two histories $h, h' \in M$ are \emph{observationally equivalent},
written $h \sim h'$, if $\Pi(h) = \Pi(h')$.
\end{definition}

\begin{remark}[On the category of $M$]
The appropriate category for $M$ varies by application. For quantum field
theory, $M$ is a space of distributional field configurations with the
weak-* topology. For the Krein-space formulation of higher-derivative
gravity, $M$ is a Krein space (an indefinite inner-product space with
a $\mathcal{K}_+ \oplus \mathcal{K}_-$ decomposition). For classical
mechanics, $M$ is a cotangent bundle. For cognitive models, $M$ is a
space of neural activation histories. For institutional models, $M$ is
a space of decision trajectories. The framework requires only that $M$
be measurable and that $\Pi$ be surjective and measurable. Additional
structure --- topology, metric, linear structure --- may be present and
will be used when available, but is not presupposed.
\end{remark}

\subsection{Admissibility and the Projection Sufficiency Principle}

In applications, not every history in $M$ is physically, cognitively, or
institutionally realizable. An admissibility predicate selects the
realizable ones.

\begin{definition}[Admissibility]
Given a projection setting $(M, O, \Pi)$, an \emph{admissibility predicate}
is a measurable function $A_O : O \to \{0,1\}$. The induced predicate
on $M$ is $A_M(h) = A_O(\Pi(h))$. A history $h$ is \emph{admissible} if
$A_M(h) = 1$.
\end{definition}

The definition places admissibility on $O$, not on $M$. This is the
central methodological commitment of the framework.

\begin{theorem}[Projection Sufficiency]
\label{thm:proj_suff}
If physical, cognitive, or institutional predictions depend only on
$A_O(\Pi(h))$, then no condition on $h \in M$ is necessary unless it
changes $\Pi(h)$.
\end{theorem}

\begin{proof}
Let $h, h' \in M$ with $\Pi(h) = \Pi(h')$. Then
\[A_M(h) = A_O(\Pi(h)) = A_O(\Pi(h')) = A_M(h').\] Any property that distinguishes $h$ from $h'$
while leaving $\Pi(h)$ unchanged has no effect on admissibility. Since
all predictions are functions of admissibility and observations, and since
both are determined by $\Pi(h)$, no condition on fiber structure can affect
any prediction.
\end{proof}

\begin{corollary}[The Injectivity Special Case]
\label{cor:injective}
When $\Pi$ is injective, $\Pi(h) = \Pi(h')$ implies $h = h'$, so every
fiber is a singleton. There are no hidden directions in the fiber, and
every internal property of $h$ is directly observable. This is the
condition corresponding to the Hilbert Assumption in quantum mechanics
(the identification of the state space with the observable space),
naïve realism in epistemology (the identification of mental content with
external reality), and the legibility assumption in institutional theory
(the identification of an organization's behavior with its internal
dynamics).
\end{corollary}

\begin{remark}[Pathologies in $M$ versus pathologies in $O$]
\label{rem:pathologies}
A \emph{coordinate pathology} is a property of some $h \in M$ that has
no bearing on $\Pi(h)$. An \emph{observable pathology} is a violation of
admissibility in $O$: $A_O(\Pi(h)) = 0$. By Theorem~\ref{thm:proj_suff},
coordinate pathologies have no observable consequences. Whether a
specific pathology in $M$ induces an observable pathology depends on
the structure of $\Pi$, not on the pathology alone. This elementary
observation underlies the resolution of the ghost problem in higher-derivative
quantum gravity, the resolution of the Ostrogradsky instability in
gravitational cosmology, and the analysis of institutional dysfunction
in large organizations that continue to function despite severe internal
contradictions.
\end{remark}


%======================================================================
\section{Reconstruction Defect and Irreversibility}
%======================================================================

\subsection{The Reconstruction Problem}

Given an observation $o = \Pi(h)$, can the generating history $h$ be
recovered? This is the reconstruction problem. When $\Pi$ is injective,
the answer is yes: the unique preimage $\Pi^{-1}(o)$ recovers $h$
exactly. When $\Pi$ is many-to-one, the fiber $F_o$ contains multiple
histories, and reconstruction must choose among them.

\begin{definition}[Reconstruction Map and Defect]
\label{def:reconstruction}
A \emph{reconstruction map} is a measurable function
$R : O \to M$ such that $\Pi(R(o)) = o$ for all $o \in O$ (that is,
$R$ is a section of $\Pi$). Given a metric $d$ on $M$, the
\emph{reconstruction defect} of $R$ at $h$ is
\begin{equation}
  \Delta_R(h) = d\bigl(h,\, R(\Pi(h))\bigr).
\end{equation}
The \emph{worst-case reconstruction defect} of $R$ is
$\|\Delta_R\|_\infty = \sup_{h \in M} \Delta_R(h)$.
The \emph{minimum reconstruction defect} of the projection is
\begin{equation}
  \delta(\Pi) = \inf_R \|\Delta_R\|_\infty,
\end{equation}
where the infimum is taken over all reconstruction maps.
\end{definition}

\begin{theorem}[Irreversibility from Non-Injectivity]
\label{thm:irreversibility}
Let $\Pi : M \to O$ and let $d$ be a metric on $M$. Then
\begin{equation}
  \delta(\Pi)
  \;\geq\;
  \frac{1}{2}\sup_{o \in O}\operatorname{diam}(F_o),
  \label{eq:defect_diam}
\end{equation}
where $\operatorname{diam}(F_o) = \sup_{h,h' \in F_o} d(h,h')$. In
particular, if some fiber $F_o$ has positive diameter then every
reconstruction map $R$ has positive worst-case defect on that fiber:
$\sup_{h \in F_o} \Delta_R(h) \geq \frac{1}{2}\operatorname{diam}(F_o) > 0$.
\end{theorem}

\begin{proof}
Fix any $o \in O$ and any two points $h, h' \in F_o$. For any reconstruction
map $R$, the triangle inequality gives
\begin{equation*}
  d(h,h') \leq d(h, R(o)) + d(R(o), h')
  = \Delta_R(h) + \Delta_R(h').
\end{equation*}
Therefore $\max(\Delta_R(h), \Delta_R(h')) \geq \frac{1}{2}d(h,h')$.
Taking the supremum over $h, h' \in F_o$ and then over all $o$ and all
reconstruction maps $R$:
\begin{equation*}
  \delta(\Pi) = \inf_R \sup_{h \in M} \Delta_R(h)
  \geq \frac{1}{2}\sup_{o \in O}\sup_{h,h' \in F_o} d(h,h')
  = \frac{1}{2}\sup_{o \in O}\operatorname{diam}(F_o).
\end{equation*}
When $d$ separates points, any non-singleton fiber has positive diameter,
so $\delta(\Pi) > 0$ whenever $\Pi$ is non-injective.
\end{proof}

\begin{remark}
Theorem~\ref{thm:irreversibility} is the precise mathematical content of
the claim that information is lost under projection. The reconstruction
defect $\delta(\Pi)$ is a measure of how much is lost: it is the minimum
error that any reconstruction must make, regardless of how cleverly it
tries to infer the generating history from the observation. Irreversibility
is not a property of the dynamics of the system but of the projection by
which the system is observed. A dynamical system whose intrinsic evolution
is time-reversible appears irreversible to any observer whose projection
$\Pi$ is non-injective.
\end{remark}

\subsection{The Defect Spectrum}

The minimum worst-case defect $\delta(\Pi)$ is a single number. It is
useful to have a finer characterisation of how reconstruction defect is
distributed across $M$.

\begin{definition}[Defect Spectrum]
Given $(M, O, \Pi, d)$, the \emph{defect spectrum} is the function
$\delta_* : [0,\infty) \to [0,1]$ defined by
\begin{equation}
  \delta_*(t) = \sup_R\, \mu\!\bigl(\{h \in M : \Delta_R(h) \leq t\}\bigr),
\end{equation}
where $\mu$ is a reference probability measure on $M$ and the supremum
is over all reconstruction maps $R : O \to M$. Thus $\delta_*(t)$ is
the largest fraction of histories that the best reconstruction map can
recover to within distance $t$.
\end{definition}

The defect spectrum interpolates between $\delta_*(0) = 0$ (when $\Pi$
is non-injective, no reconstruction map achieves zero defect everywhere)
and $\delta_*(\infty) = 1$ (the entire history space is recovered to
within infinite tolerance). Its shape characterises the nature of the
information loss:

\begin{proposition}[Defect Spectrum at Zero]
$\delta_*(0) = 1$ if and only if $\Pi$ is injective. When $\Pi$ is
non-injective, $\delta_*(0) < 1$: even the best reconstruction map fails
on a positive-measure set of histories.
\end{proposition}

\begin{proof}
If $\Pi$ is injective, $R = \Pi^{-1}$ satisfies $\Delta_R(h) = 0$ for all
$h$, so $\delta_*(0) = \mu(M) = 1$. If $\Pi$ is non-injective, some fiber
$F_o$ is non-singleton; by Theorem~\ref{thm:irreversibility}, any $R$
incurs positive defect on at least one element of $F_o$, so
$\mu(\{h : \Delta_R(h) = 0\}) < \mu(M) = 1$ for every $R$, giving
$\delta_*(0) < 1$.
\end{proof}

\subsection{Residual Uncertainty and Reconstruction Defect}

The information lost under projection is quantified by the
\emph{conditional entropy of $h$ given $\Pi(h)$}:
\begin{equation}
  H(h \mid \Pi(h)) = \sum_{o \in O} \nu(o)\, H(\mu_{F_o}),
  \label{eq:cond_entropy}
\end{equation}
where $\mu_{F_o}$ is the conditional distribution of $h$ on the fiber
$F_o$ and $\nu = \Pi_*\mu$. This is the residual uncertainty about
the generating history that survives after observing $\Pi(h) = o$:
the average entropy within fibers, not the entropy of the projected
distribution.

\begin{remark}[Why $H(\Pi_*\mu)$ is the wrong bound]
The Shannon entropy $H(\nu) = H(\Pi_*\mu)$ measures uncertainty about
which fiber an observation falls in; it does not measure uncertainty about
the history \emph{within} a fiber. A projection that collapses everything
into a single fiber has $H(\nu) = 0$ but $H(h \mid \Pi(h))$ arbitrarily
large. The operative quantity for reconstruction error is the conditional
entropy $H(h \mid \Pi(h))$, not the marginal $H(\Pi_*\mu)$.
\end{remark}

\begin{theorem}[Fano-Type Bound on Reconstruction]
\label{thm:entropy_bound}
Let $(M, \mathcal{F}, \mu)$ be a probability space with $M$ finite, and
let $R : O \to M$ be any reconstruction map. Then
\begin{equation}
  P_e(R) \;\geq\; \frac{H(h \mid \Pi(h)) - \log 2}{\log|M|},
  \label{eq:fano}
\end{equation}
where $P_e(R) = \mu(\{h : R(\Pi(h)) \neq h\})$ is the error probability.
Consequently, whenever $H(h \mid \Pi(h)) > \log 2$, no reconstruction map
achieves zero error probability.
\end{theorem}

\begin{proof}
Apply Fano's inequality \cite{cover2006} to the channel $h \to \Pi(h)
\to R(\Pi(h))$. With input $h$, output $R(\Pi(h))$, and error event
$\{R(\Pi(h)) \neq h\}$, Fano's inequality gives
$H(h \mid \Pi(h)) \leq H(P_e) + P_e\log(|M|-1) \leq \log 2 + P_e\log|M|$,
which rearranges to~(\ref{eq:fano}).
\end{proof}

\begin{remark}[Metric Defect via Rate-Distortion]
For a metric lower bound on $\mathbb{E}[\Delta_R]$ rather than on error
probability, rate-distortion theory gives
$\mathbb{E}_\mu[\Delta_R] \geq D^*(H(h \mid \Pi(h)))$,
where $D^*(I)$ is the rate-distortion function of the fiber-conditional
distribution at information constraint $I$ \cite{cover2006}. In both
cases the operative quantity is $H(h \mid \Pi(h))$~(\ref{eq:cond_entropy}).
\end{remark}

\begin{corollary}[Thermodynamic Interpretation]
\label{cor:thermo}
When $M$ is the microstate space of a thermodynamic system, $O$ is the
macrostate space, and $\Pi$ is the standard coarse-graining map, the
Boltzmann entropy $S = k_B \log |F_o|$ (for uniform measure on $M$) is
a measure of reconstruction defect: the logarithm of the number of
microstates compatible with the macrostate is the log-volume of the fiber.
The second law of thermodynamics, under this reading, is the statement
that a system initialized in a small fiber and evolving under a
measure-preserving dynamics will, under coarse-graining, appear to move
toward states with larger fibers, increasing the reconstruction defect.
\end{corollary}

\subsection{Irreversibility as Accumulated Defect}

For systems with a dynamical evolution, irreversibility is not merely the
existence of reconstruction defect but its persistence and growth under
the dynamics.

\begin{definition}[Dynamical Irreversibility]
Let $\phi_t : M \to M$ be a flow on $M$ (the dynamics), and let $\Pi$
be the observation projection. The \emph{accumulated reconstruction
defect} at time $T$ starting from $h$ is
\begin{equation}
  \mathcal{I}(h, T) = \inf_R \int_0^T \Delta_R(\phi_t(h))\, dt.
\end{equation}
A system is \emph{dynamically irreversible at $h$} if $\mathcal{I}(h,T) > 0$
for all $T > 0$ and all reconstruction maps $R$.
\end{definition}

\begin{proposition}[Conditions for Dynamical Irreversibility]
If $\Pi \circ \phi_t$ is non-injective for $\mu$-almost every $t > 0$,
then the system is dynamically irreversible at $\mu$-almost every $h \in M$.
\end{proposition}

\begin{proof}
Non-injectivity of $\Pi \circ \phi_t$ means there exist $h \neq h'$ with
$\Pi(\phi_t(h)) = \Pi(\phi_t(h'))$. By Theorem~\ref{thm:irreversibility},
any reconstruction map incurs positive defect at either $\phi_t(h)$ or
$\phi_t(h')$ for each such $t$. Integrating over $t$ gives positive
accumulated defect.
\end{proof}

\begin{remark}[The Arrow of Time]
The arrow of time, in the framework of reconstruction defect, is the
direction in which $\mathcal{I}(h,T)$ grows. It is not a property of
the fundamental dynamics $\phi_t$ --- which may be time-reversible at
the level of $M$ --- but of the projection $\Pi$ under which the system
is observed. A perfectly reversible dynamics observed through a
non-injective projection appears irreversible because the accumulated
defect grows with $T$: the observer loses the ability to determine which
of the histories compatible with the current observation was the actual
generating history. This is the precise sense in which the arrow of time
is a projection artifact.
\end{remark}


%======================================================================
\section{The Distinction Functional}
%======================================================================

\subsection{The Problem of Comparing Interfaces}

Two observers equipped with different projections $\Pi_1$ and $\Pi_2$
will assign different probability distributions to the same underlying
space of histories. How should the difference between these distributions
be measured?

The answer must satisfy several constraints. First, it must be
\emph{interface-native}: the comparison should require no reference to
the underlying space $M$ or the projections $\Pi_1, \Pi_2$. Once the
distributions $P = \Pi_{1*}\mu_1$ and $Q = \Pi_{2*}\mu_2$ on $O$ are
given, the comparison should proceed entirely within $O$.

Second, the measure should be \emph{symmetric}: the distance from $P$
to $Q$ should equal the distance from $Q$ to $P$. There is no
privileged interface.

Third, the measure should be \emph{operationally meaningful}: it should
have an interpretation in terms of distinguishability, information, or
some quantity that can in principle be estimated from observations.

Fourth, it should be \emph{a metric}: zero if and only if $P = Q$,
satisfying the triangle inequality.

The Jensen--Shannon divergence satisfies all four constraints. We develop
its properties systematically in this section.

\subsection{Definition and Basic Properties}

\begin{definition}[Jensen--Shannon Divergence and Distance]
\label{def:jsd}
Let $P$ and $Q$ be probability distributions on a measurable space $O$.
Let $M_{PQ} = \frac{1}{2}(P + Q)$ be their equal-weight mixture. The
\emph{Jensen--Shannon divergence} is
\begin{equation}
  \operatorname{JSD}(P \| Q)
  = \frac{1}{2} D_{\mathrm{KL}}(P \| M_{PQ})
    + \frac{1}{2} D_{\mathrm{KL}}(Q \| M_{PQ}),
\end{equation}
where $D_{\mathrm{KL}}(P \| Q) = \sum_x P(x) \log(P(x)/Q(x))$ is the
Kullback--Leibler divergence. Equivalently,
\begin{equation}
  \operatorname{JSD}(P \| Q)
  = H(M_{PQ}) - \frac{1}{2}\bigl(H(P) + H(Q)\bigr),
  \label{eq:jsd_entropy}
\end{equation}
where $H(\cdot)$ is the Shannon entropy. The \emph{Jensen--Shannon
distance} is $d_{\mathrm{int}}(P,Q) = \operatorname{JSD}^{1/2}(P \| Q)$.
\end{definition}

\begin{proposition}[Metric Properties]
\label{prop:metric}
The Jensen--Shannon distance $d_{\mathrm{int}}$ is a metric on the space
of probability distributions on $O$. Specifically:
\begin{enumerate}[label=(\roman*)]
  \item $d_{\mathrm{int}}(P,Q) \geq 0$, with equality if and only if $P = Q$;
  \item $d_{\mathrm{int}}(P,Q) = d_{\mathrm{int}}(Q,P)$;
  \item $d_{\mathrm{int}}(P,R) \leq d_{\mathrm{int}}(P,Q) + d_{\mathrm{int}}(Q,R)$.
\end{enumerate}
\end{proposition}

\begin{proof}
Properties (i) and (ii) follow immediately from the definition.
$\operatorname{JSD}(P\|Q) = 0$ iff both KL terms vanish, iff $P = M_{PQ} = Q$.
Symmetry is manifest from the equal-weight mixture. The triangle
inequality for $d_{\mathrm{int}} = \operatorname{JSD}^{1/2}$ was established
by Endres and Schindelin \cite{endres2003} using the concavity of
$\sqrt{\cdot}$ and the properties of KL divergence. The key step is that
$\operatorname{JSD}^{1/2}$ satisfies the triangle inequality while
$\operatorname{JSD}$ alone does not.
\end{proof}

\begin{remark}
The need to take the square root --- that $\operatorname{JSD}^{1/2}$ is
a metric while $\operatorname{JSD}$ is not --- is the information-geometric
signature of the curvature we will compute in Section~\ref{sec:geometry}.
In a flat information geometry, the square root would be unnecessary; the
need for it indicates that the space of probability distributions is
positively curved in the Fisher--Rao metric.
\end{remark}

\subsection{The Distinction Functional}

\begin{definition}[Distinction Functional]
\label{def:distinction}
The \emph{distinction functional} between observable interfaces $P$ and
$Q$ is
\begin{equation}
  \mathcal{D}(P,Q)
  = H(M_{PQ}) - \frac{1}{2}\bigl(H(P) + H(Q)\bigr)
  = \operatorname{JSD}(P \| Q).
\end{equation}
\end{definition}

The entropy decomposition~(\ref{eq:jsd_entropy}) gives each term a
precise operational meaning, and their combination characterises the
information carried by the distinction between $P$ and $Q$.

\begin{theorem}[Operational Meaning of the Distinction Functional]
\label{thm:distinction_meaning}
The distinction functional $\mathcal{D}(P,Q)$ measures the expected
reduction in uncertainty about which interface generated an observation,
given that the observation has been made. Equivalently, it is the
information that is destroyed when the distinction between $P$ and $Q$
is forgotten by replacing both with their mixture $M_{PQ}$.
\end{theorem}

\begin{proof}
Consider an observer who knows an observation $x \in O$ was generated by
one of two sources, each equally likely: source 1 using distribution $P$,
source 2 using distribution $Q$. The prior uncertainty about the source is
$H(\text{source}) = \log 2$. The observation $x$ is drawn from the mixture
$M_{PQ}$. After seeing $x$, the posterior uncertainty about the source is
$H(\text{source} \mid x) = -\sum_{s \in \{1,2\}} P(s \mid x) \log P(s\mid x)$.
The expected posterior uncertainty over observations is
$\sum_x M_{PQ}(x) H(\text{source} \mid x)$.

By the chain rule of entropy applied to the joint $(x, \text{source})$:
\begin{align*}
  H(x, \text{source})
  &= H(\text{source}) + H(x \mid \text{source})
   = \log 2 + \frac{1}{2}(H(P) + H(Q)), \\
  H(x, \text{source})
  &= H(x) + H(\text{source} \mid x)
   = H(M_{PQ}) + H(\text{source} \mid x).
\end{align*}
Therefore $H(\text{source} \mid x) = \log 2 + \frac{1}{2}(H(P)+H(Q)) - H(M_{PQ})
= \log 2 - \mathcal{D}(P,Q)$. The information gain from observing $x$ is
$\log 2 - H(\text{source}\mid x) = \mathcal{D}(P,Q)$. When the distinction
is forgotten ($P$ and $Q$ are merged into $M_{PQ}$), this information
gain vanishes: the expected uncertainty about the source becomes $\log 2$
(no information about the source), and the distinction functional equals
exactly the information that was lost.
\end{proof}

\begin{corollary}[Monotonicity and Data Processing]
\label{cor:data_proc}
The distinction functional satisfies the data processing inequality: for
any measurable map $f : O \to O'$,
\begin{equation}
  \mathcal{D}(f_*P,\, f_*Q) \leq \mathcal{D}(P,Q).
\end{equation}
Any further processing of observations can only reduce the distinction
between the interfaces.
\end{corollary}

\begin{proof}
This follows from the data processing inequality for KL divergence:
\[D_{\mathrm{KL}}(f_*P \| f_*M_{PQ}) \leq D_{\mathrm{KL}}(P \| M_{PQ}),\]
and similarly for $Q$.
\end{proof}

\subsection{Interface-Nativity}

\begin{proposition}[Interface-Nativity]
\label{prop:interface_native}
The distinction functional $\mathcal{D}(P,Q)$ is interface-native: its
definition and computation require no reference to any underlying space
$M$ or projection $\Pi$.
\end{proposition}

\begin{proof}
Definition~\ref{def:distinction} uses only the distributions $P$ and $Q$
on the observable space $O$, their mixture $M_{PQ}$, and the Shannon
entropy $H$. No element of $M$ appears, no projection $\Pi$ is referenced,
and no admissibility manifold is invoked. The computation can be performed
by any agent who has access to $O$ and can sample from $P$ and $Q$,
without knowledge of how those distributions were generated.
\end{proof}

\begin{remark}[Contrast with the Inf-Sup Formulation]
An alternative distance between theories with underlying spaces
$M^{(1)}$ and $M^{(2)}$ would be
\begin{equation}
  d_{\mathrm{ext}}(M^{(1)}, M^{(2)})
  = \inf_\phi \sup_{h} \bigl\|\Pi_1(h) - \Pi_2(\phi(h))\bigr\|_O,
\end{equation}
where the infimum runs over admissibility-preserving maps $\phi$.
This formulation is not interface-native: the infimum over $\phi$
requires access to $M^{(1)}$ and $M^{(2)}$. It frames the question as
``how similar can the internal structures be made to look?'' rather than
``how similar are the observable outputs?'' The distinction functional
answers the second question without invoking the first, and is therefore
the appropriate primitive for an interface-native theory.
\end{remark}

\subsection{Gauge Orbits and Hidden Directions}

\begin{theorem}[Gauge Orbits as Zero-Distinction Fibers]
\label{thm:gauge}
Let $(M, O, \Pi)$ be a projection setting. The observational equivalence
relation $h \sim h'$ (Definition~1.1) is the zero-level set of the
induced distinction functional between the Dirac measures at the projected
points:
\begin{equation}
  h \sim h' \quad\iff\quad
  \mathcal{D}(\delta_{\Pi(h)},\, \delta_{\Pi(h')}) = 0,
\end{equation}
where $\delta_o$ denotes the Dirac measure concentrated at $o \in O$.
\end{theorem}

\begin{proof}
The distinction functional compares probability distributions on $O$. For
point masses, $\mathcal{D}(\delta_o, \delta_{o'}) = \operatorname{JSD}(\delta_o
\| \delta_{o'})$. Since $H(\delta_o) = 0$ for any point mass and
$H(\frac{1}{2}(\delta_o + \delta_{o'})) = \log 2$ when $o \neq o'$ and $0$
when $o = o'$, we get $\mathcal{D}(\delta_o, \delta_{o'}) = \log 2$ when
$o \neq o'$ and $0$ when $o = o'$. Therefore
\[\mathcal{D}(\delta_{\Pi(h)},\, \delta_{\Pi(h')}) = 0
  \;\iff\; \Pi(h) = \Pi(h'),\]
which is the definition of $h \sim h'$. Gauge orbits are the fibers
$F_o = \Pi^{-1}(o)$: the zero-distinction equivalence classes.
\end{proof}

\begin{definition}[Hidden Admissibility Direction]
\label{def:hidden}
Let $h_\epsilon$ be a smooth curve in $M$ with $h_0 = h$. The
\emph{first-order observable variation} is
\begin{equation}
  \delta_{\mathrm{obs}}(h, \dot h)
  = \left.\frac{d}{d\epsilon}
    d_{\mathrm{int}}\!\bigl(\Pi(h),\, \Pi(h_\epsilon)\bigr)
    \right|_{\epsilon=0}.
\end{equation}
The direction $\dot h = \frac{dh_\epsilon}{d\epsilon}|_{\epsilon=0}$
is \emph{hidden to first order} if $\delta_{\mathrm{obs}}(h, \dot h) = 0$.
It is \emph{fully hidden} if $\mathcal{D}(\delta_{\Pi(h)}, \delta_{\Pi(h_\epsilon)}) = 0$
for all sufficiently small $\epsilon > 0$.
\end{definition}

\begin{example}[Ghost States as Hidden Directions]
\label{ex:ghosts}
In the Krein-space formulation of higher-derivative quantum gravity,
$M = \mathcal{K} = \mathcal{K}_+ \oplus \mathcal{K}_-$ and $\Pi$ maps
each state to its observable transition probabilities via
$P(F \mid I) = \mathrm{tr}_+((\Pi_F S \Pi_I)^\dagger(\Pi_F S \Pi_I))$.
A ghost state $g \in \mathcal{K}_-$ generates a deformation $h_\epsilon = h
+ \epsilon g$ in $M$. If ghost-parity symmetry holds ($[G, S] = 0$), then
the contribution of $g$ to any observable transition probability cancels
in the trace over $\mathcal{K}$, so $\Pi(h + \epsilon g) = \Pi(h)$ for
all $\epsilon$. The deformation in the $g$-direction is fully hidden: the
ghost is a hidden admissibility direction with $\mathcal{D}(\Pi(h),
\Pi(h + \epsilon g)) = 0$ for all $\epsilon$. Whether ghost-parity symmetry
holds in the full spin-2 sector of quadratic gravity remains open; the
present framework identifies this as precisely the question of whether the
ghost direction is fully hidden.
\end{example}

\begin{example}[Gauge Transformations as Hidden Directions]
In Yang--Mills theory, $M$ is the space of gauge field configurations
$A_\mu$ and $\Pi$ is the map to gauge-invariant observables (Wilson loops,
field strength curvature, etc.). A gauge transformation $A_\mu \mapsto
A_\mu + D_\mu\lambda$ is a deformation along a hidden direction: by gauge
invariance, $\Pi(A_\mu + D_\mu\lambda) = \Pi(A_\mu)$, so the deformation
is fully hidden. The gauge orbit through $A_\mu$ is the fiber $F_{\Pi(A_\mu)}$.
\end{example}

\begin{example}[Cognitive Hidden Directions]
\label{ex:cognitive}
For a cognitive agent with internal representation space $M$ and
perceptual projection $\Pi$ to sensory observations, a hidden direction
is an internal change that leaves all perceptual outputs unchanged. This
includes: changes in sub-threshold neural activity, shifts in attentional
allocation that do not affect behavioral output, alterations in implicit
beliefs that have not yet influenced action, and unconscious emotional
states that are not expressed in observable behavior. The distinction
functional $\mathcal{D}(\delta_{\Pi(h)}, \delta_{\Pi(h+\epsilon\delta h)})$ measures
whether the internal change $\delta h$ is cognitively legible to an
external observer.
\end{example}


%======================================================================
\section{The Geometry of Interface Space}
\label{sec:geometry}
%======================================================================

\subsection{The Space of Theories as a Metric Space}

Let $\mathcal{M}$ denote the space of all observable interfaces: all
probability distributions that can be generated by some admissible
projection setting $(M, O, \Pi, \mu)$. Two interfaces $P, Q \in \mathcal{M}$
are compared by the distinction functional $\mathcal{D}(P,Q) =
\operatorname{JSD}(P\|Q)$, and the interface distance $d_{\mathrm{int}} =
\mathcal{D}^{1/2}$ makes $\mathcal{M}$ a metric space by
Proposition~\ref{prop:metric}.

\begin{definition}[Theory Space]
\label{def:theory_space}
The \emph{theory space} of a given observable domain $O$ is the metric
space $(\mathcal{M}, d_{\mathrm{int}})$ of probability distributions on
$O$, equipped with the Jensen--Shannon distance.
\end{definition}

\begin{remark}
The elements of $\mathcal{M}$ are observable distributions, not internal
models. Two theories with very different internal structures --- different
$M$'s, different $\Pi$'s, different dynamics on $M$ --- may occupy the
same point in $(\mathcal{M}, d_{\mathrm{int}})$ if their projected
distributions agree. The theory space is not a space of theories in the
traditional sense (axiomatic systems, Lagrangians, Hamiltonians) but a
space of what theories produce: observable distributions over outcomes.
It is the appropriate space for empirical comparison.
\end{remark}

\subsection{Local Riemannian Structure: The Fisher Metric}

For a parametric family of distributions $\{P_\lambda\}_{\lambda \in \Lambda}$,
the distinction functional induces a Riemannian metric on $\Lambda$.

\begin{theorem}[Fisher Metric from JSD]
\label{thm:fisher}
Let $\lambda \mapsto P_\lambda$ be a smooth map from an open subset
$\Lambda \subseteq \mathbb{R}^n$ to probability distributions on $O$.
Then the Jensen--Shannon divergence satisfies
\begin{equation}
  \operatorname{JSD}(P_\lambda,\, P_{\lambda+\epsilon})
  = \frac{1}{8}\, g^F_{ij}(\lambda)\,\epsilon^i\epsilon^j
    + O(|\epsilon|^3),
  \label{eq:jsd_fisher}
\end{equation}
so that $d_{\mathrm{int}}^2 = \operatorname{JSD} = \frac{1}{8}g^F_{ij}\epsilon^i\epsilon^j + O(|\epsilon|^3)$,
where
\begin{equation}
  g^F_{ij}(\lambda)
  = \sum_{x \in O} P_\lambda(x)\,
    \frac{\partial \log P_\lambda(x)}{\partial \lambda^i}\,
    \frac{\partial \log P_\lambda(x)}{\partial \lambda^j}
\end{equation}
is the Fisher information metric on $\Lambda$.
\end{theorem}

\begin{proof}
Let $\delta P(x) = P_{\lambda+\epsilon}(x) - P_\lambda(x) =
\partial_i P_\lambda(x)\,\epsilon^i + O(|\epsilon|^2)$ and
$M(x) = P_\lambda(x) + \frac{1}{2}\delta P(x)$. Expanding
$D_{\mathrm{KL}}(P_\lambda \| M)$:
\begin{equation*}
  D_{\mathrm{KL}}(P_\lambda \| M)
  = \sum_x P_\lambda(x)\log\frac{P_\lambda(x)}{M(x)}
  = \sum_x P_\lambda(x)\log\!\left(1 - \frac{\delta P(x)}{2M(x)}\right).
\end{equation*}
Since $M(x) = P_\lambda(x)(1 + O(|\epsilon|))$, we have
$\frac{\delta P(x)}{2M(x)} = \frac{\delta P(x)}{2P_\lambda(x)} + O(|\epsilon|^2)$.
Using $\log(1-u) = -u - u^2/2 + O(u^3)$:
\begin{equation*}
  D_{\mathrm{KL}}(P_\lambda \| M)
  = \frac{1}{8}\sum_x \frac{(\delta P(x))^2}{P_\lambda(x)} + O(|\epsilon|^3).
\end{equation*}
By symmetry $D_{\mathrm{KL}}(P_{\lambda+\epsilon}\|M)$ contributes the
same leading term. Combining both KL terms
($\operatorname{JSD} = \frac{1}{2}D_{\mathrm{KL}}(P_\lambda\|M) +
\frac{1}{2}D_{\mathrm{KL}}(P_{\lambda+\epsilon}\|M)$), each contributing
$\frac{1}{8}\sum_x (\delta P)^2/P_\lambda$ at leading order:
\begin{equation*}
  \operatorname{JSD}(P_\lambda,\, P_{\lambda+\epsilon})
  = \frac{1}{8}\sum_x \frac{(\delta P(x))^2}{P_\lambda(x)} + O(|\epsilon|^3)
  = \frac{1}{8}\, g^F_{ij}(\lambda)\,\epsilon^i\epsilon^j + O(|\epsilon|^3),
\end{equation*}
where $g^F_{ij} = \sum_x P_\lambda(x)(\partial_i\log P_\lambda)
(\partial_j\log P_\lambda)$ is the Fisher information matrix. Since
$d_{\mathrm{int}} = \operatorname{JSD}^{1/2}$, we have
$d_{\mathrm{int}}^2 = \operatorname{JSD} = \frac{1}{8}g^F_{ij}\epsilon^i
\epsilon^j + O(|\epsilon|^3)$, confirming~(\ref{eq:jsd_fisher}).
\end{proof}

\begin{remark}[Relationship to Information Geometry]
Theorem~\ref{thm:fisher} establishes that $(\mathcal{M}, d_{\mathrm{int}})$
is locally a statistical manifold \cite{amari2000} with the Fisher--Rao
metric. The factor $\frac{1}{8}$ is a convention-dependent rescaling
(it becomes $\frac{1}{2}$ in the natural units of information geometry
where logarithms are taken base $e$). The key point is that the
information geometry of the space of theories is determined entirely by
the Fisher information of the observable distributions, without reference
to the internal models that generate them.
\end{remark}

\subsection{Curvature and Its Interpretation}

The Riemannian curvature of $(\mathcal{M}, g^F)$ is the standard
Riemannian curvature of the statistical manifold with Fisher metric. It
has been computed explicitly for many exponential families \cite{amari2000}.
Here we give its operational interpretation in the theory-comparison setting.

\begin{definition}[Sectional Curvature of Theory Space]
The \emph{sectional curvature} $\kappa(\sigma)$ of $(\mathcal{M}, g^F)$
at a point $P_\lambda$ in a 2-plane $\sigma \subset T_{P_\lambda}\mathcal{M}$
is the Gaussian curvature of the 2-dimensional surface swept out by
geodesics from $P_\lambda$ in directions $\sigma$.
\end{definition}

\begin{theorem}[Operational Interpretation of Curvature]
\label{thm:curvature}
At a theory $P_\lambda \in \mathcal{M}$:
\begin{enumerate}[label=(\roman*)]
  \item Positive sectional curvature in direction $\sigma$ means that
        theories near $P_\lambda$ in the $\sigma$-direction are more
        similar in their observable predictions than their parameter-space
        distance suggests. Geodesic spheres of radius $r$ in theory space
        have \emph{smaller} volume than Euclidean spheres of radius $r$
        (the Jacobi fields converge), so more parameter-space directions
        map into the same observable neighbourhood. This is the geometry
        of \emph{overfitting}: many internal models produce
        observationally indistinguishable predictions.
  \item Negative sectional curvature in direction $\sigma$ means that
        theories near $P_\lambda$ in the $\sigma$-direction diverge in
        their observable predictions faster than their parameter-space
        distance suggests. Geodesic spheres of radius $r$ have \emph{larger}
        volume than Euclidean spheres of radius $r$ (the Jacobi fields
        diverge), so small parameter changes produce large observable
        changes. This is the geometry of \emph{high falsifiability}:
        small changes in the theory produce large changes in predictions.
  \item Zero sectional curvature is the Euclidean regime:
        parameter-space distance and observable-prediction distance
        scale proportionally, and the theory is locally well-identified.
\end{enumerate}
\end{theorem}

\begin{proof}[Proof of (i) and (ii)]
The Jacobi equation governs the evolution of nearby geodesics: if
$J(t)$ is a Jacobi field along a unit-speed geodesic $\gamma(t)$, then
$\|J(t)\| \approx t - \frac{\kappa}{6}t^3 + O(t^5)$ for small $t$, where
$\kappa$ is the sectional curvature in the plane $\{\dot\gamma, J(0)\}$.
When $\kappa > 0$, $\|J(t)\| < t$: nearby geodesics are closer than
linearly predicted, meaning theories at parameter distance $t$ are
$d_{\mathrm{int}}$-closer than expected, i.e., more observationally
similar. When $\kappa < 0$, $\|J(t)\| > t$: theories at parameter
distance $t$ are more observationally different than the Euclidean
approximation predicts, i.e., more falsifiable.
\end{proof}

\begin{example}[Curvature in Specific Families]
For exponential family distributions $P_\lambda(x) = \exp(\lambda \cdot
T(x) - A(\lambda))$ with sufficient statistic $T$, the Fisher metric is
$g^F_{ij} = \partial_i\partial_j A(\lambda)$ (the Hessian of the
log-partition function). The curvature is determined by third-order
cumulants of $T$ under $P_\lambda$. In the Gaussian family
$\mathcal{N}(\mu, \sigma^2)$, the sectional curvature in the
$(\mu, \log\sigma)$ plane is $-\frac{1}{2}$: the space of Gaussian
distributions is a hyperbolic space. This means distinguishing Gaussian
theories by their observable predictions is easier than parameter-space
distance would suggest --- small changes in $(\mu, \sigma)$ produce large
changes in $d_{\mathrm{int}}$.
\end{example}

\subsection{Geodesics as Paths of Minimal Accumulated Distinction}

\begin{definition}[Geodesic in Theory Space]
A \emph{geodesic} in $(\mathcal{M}, g^F)$ is a smooth curve
$\lambda(t)$ in the theory parameter space that satisfies the
geodesic equation $\nabla_{\dot\lambda}\dot\lambda = 0$ with respect to
the Levi-Civita connection of the Fisher metric. Its interpretation:
the geodesic between theories $P$ and $Q$ is the path through theory
space that minimises the total accumulated distinction between successive
theories along the path,
\begin{equation}
  \int_0^1 d_{\mathrm{int}}\!\bigl(P_{\lambda(t)},
  P_{\lambda(t+dt)}\bigr) = \int_0^1 \!\sqrt{g^F_{ij}\dot\lambda^i
  \dot\lambda^j}\, dt.
\end{equation}
\end{definition}

\begin{remark}[Scientific and Cognitive Interpretation]
The geodesic distance $d_{\mathrm{int}}(P,Q)$ between two theories is
the minimum total accumulated distinction along any path connecting them.
In the context of scientific theory change, a paradigm shift from theory
$P$ to theory $Q$ that follows the geodesic is the shift that introduces
the minimum total observational difference at each step. A shift that
follows a longer path introduces unnecessary intermediate distinctions ---
it makes claims that cannot be verified by any available observations.
The geodesic is the \emph{most parsimonious} path between two empirical
positions. Occam's razor, in this framework, is the injunction to travel
geodesically.
\end{remark}

\subsection{The Derived Interface Conjecture in Metric Terms}

The physical application of theory space geometry provides a precise
formulation of what it means for two theories of gravity to make the
same observable predictions below the Planck scale.

\begin{definition}[Sub-Planckian Interface Equivalence]
Two quantum gravity theories $T_1$ and $T_2$ are \emph{sub-Planckian
interface equivalent} if $d_{\mathrm{int}}(P_1^{(E)}, P_2^{(E)}) = 0$
for all observable energies $E \ll M_{\mathrm{Pl}}$, where $P_i^{(E)}$
is the distribution of scattering outcomes at energy $E$ predicted by
$T_i$.
\end{definition}

\begin{principle}[Derived Interface Conjecture, metric form]
Hilbert-space quantum gravity and Krein-space quadratic gravity are
sub-Planckian interface equivalent: they satisfy $d_{\mathrm{int}} = 0$
at all energies accessible to current and foreseeable experiment. Above
the Planck scale, they diverge: $d_{\mathrm{int}} > 0$, with the distance
controlled by the mass of the spin-2 ghost $m_2 = M_{\mathrm{Pl}}/\sqrt{2\beta}$.
\end{principle}

The conjecture is falsifiable in the precise sense of Theorem~\ref{thm:curvature}:
it predicts that the theory space curvature in the direction separating
the two theories is controlled by the energy scale at which they diverge,
and that no sub-Planckian measurement can distinguish them.


%======================================================================
\section{Synthesis: The Three-Level Chain}
%======================================================================

The three structures developed in Sections~2--4 form a coherent chain:

\begin{equation}
  M
  \;\xrightarrow{\;\Pi\;}
  O
  \;\xrightarrow{\;\mathcal{D}\;}
  (\mathcal{M},\, d_{\mathrm{int}}).
  \label{eq:chain}
\end{equation}

The first arrow discards internal structure: what survives is the observable
interface $P = \Pi_*\mu \in \mathcal{M}$. The information lost in this step
is measured by the reconstruction defect $\delta(\Pi)$ of Section~2. The
second arrow measures distinguishability between interfaces: what survives
is a geometry on the space of theories. The information destroyed when the
distinction between two interfaces is forgotten is measured by the
distinction functional $\mathcal{D}$ of Section~3.

The central question shifts accordingly at each level:
\begin{align*}
  \text{Level 1 (internal):}\quad
  &\text{What structures exist in } M? \\
  \text{Level 2 (interface):}\quad
  &\text{What can be inferred about } M \text{ from } O? \\
  \text{Level 3 (geometric):}\quad
  &\text{What geometry do the limitations of inference} \\
  &\quad\text{impose on }\mathcal{M}?
\end{align*}

This paper concerns primarily Levels~2 and~3. Level~1 is the domain of
specific physical, cognitive, or institutional theories; the present
framework takes whatever is at Level~1 and asks what it produces and what
structure that production imposes on the space of theories.

\begin{theorem}[Unified Translation]
\label{thm:translation}
Under the chain~(\ref{eq:chain}), the following translations hold:
\begin{enumerate}[label=(\roman*)]
  \item A gauge orbit is a zero-$\mathcal{D}$ fiber of $\Pi$
        (Theorem~\ref{thm:gauge}).
  \item A hidden admissibility direction is a tangent direction to $M$
        along which $d_{\mathrm{int}}(\Pi(h), \Pi(h_\epsilon)) = 0$
        to first order (Definition~\ref{def:hidden}).
  \item Irreversibility is persistent positive reconstruction defect
        $\Delta_R > 0$ under any reconstruction map
        (Theorem~\ref{thm:irreversibility}).
  \item The Born rule is the projection $re^{i\theta} \mapsto r^2$:
        the map from complex amplitude space to observable intensity
        space that discards the phase (a hidden direction in the fiber
        of the squaring map).
  \item Theory comparison is distance in $(\mathcal{M}, d_{\mathrm{int}})$
        (Definition~\ref{def:theory_space}).
  \item The curvature of theory space is the curvature of
        distinguishability (Theorem~\ref{thm:curvature}).
\end{enumerate}
\end{theorem}

\begin{remark}[On the Ontological Question]
The chain~(\ref{eq:chain}) does not answer the question ``what exists?''
It dissolves it. The question ``what exists in $M$?'' is replaced by
``what survives projection?'' (Level~2) and ``what structure does the
pattern of survival impose on $\mathcal{M}$?'' (Level~3). Whatever exists
in $M$ --- fields, states, decision trajectories, neural activations,
institutional dynamics --- is relevant to inquiry only to the degree that
it produces distinguishable differences in $O$. The geometry of
$(\mathcal{M}, d_{\mathrm{int}})$ is the geometry of what matters
empirically, which may be a strict subset of what exists internally.

This is not instrumentalism. It is not the claim that internal structure
is unreal. It is the claim that the epistemologically primary object ---
the object about which rational inquiry can make progress --- is the
geometry of distinguishability, not the geometry of $M$. Internal
structure is a tool for generating and explaining observable distinctions,
not an object of inquiry in its own right.
\end{remark}

%----------------------------------------------------------------------

%======================================================================
\section{Projection Chains and Composed Loss}
%======================================================================

\subsection{Composing Projections}

Many systems involve not a single projection but a chain of projections:
an internal space is projected onto an intermediate space, which is then
projected again onto the final observable space. Brains project sensory
signals onto perceptual representations, which are then projected onto
behavioral outputs. Scientific instruments project physical quantities
onto measurement readings, which are then projected onto reported values.
Institutions project individual decisions onto collective actions, which
are then projected onto public outcomes.

\begin{definition}[Projection Chain]
A \emph{projection chain} of length $n$ is a sequence of projection
settings
\begin{equation}
  M_0 \xrightarrow{\;\Pi_1\;} M_1 \xrightarrow{\;\Pi_2\;} M_2
  \;\cdots\; \xrightarrow{\;\Pi_n\;} M_n = O,
\end{equation}
with composite projection $\Pi = \Pi_n \circ \cdots \circ \Pi_1 :
M_0 \to O$. The \emph{intermediate observable spaces} are
$M_1, \ldots, M_{n-1}$.
\end{definition}

\begin{theorem}[Defect Accumulates Along Chains]
\label{thm:chain_defect}
Let $\Pi = \Pi_2 \circ \Pi_1$ be a composed projection. Then for any
metrics $d_0$ on $M_0$ and $d_1$ on $M_1$:
\begin{equation}
  \delta(\Pi) \geq \delta(\Pi_1).
\end{equation}
Composing with a further projection $\Pi_2$ cannot decrease the minimum
reconstruction defect.
\end{theorem}

\begin{proof}
Let $R : O \to M_0$ be any reconstruction map for $\Pi$. Consider
the reconstruction map $R_1 = \Pi_1 \circ R$ for $\Pi_2$: it maps
$o \in O$ to $\Pi_1(R(o)) \in M_1$. For any $h \in M_0$,
\begin{equation*}
  \Delta_{R_1}(\Pi_1(h))
  = d_1(\Pi_1(h),\, R_1(\Pi_2(\Pi_1(h))))
  = d_1(\Pi_1(h),\, \Pi_1(R(\Pi(h)))).
\end{equation*}
If $d_1$ satisfies $d_1(\Pi_1(h), \Pi_1(h')) \leq d_0(h, h')$
(i.e., $\Pi_1$ is 1-Lipschitz), then
\begin{equation*}
  \Delta_{R_1}(\Pi_1(h)) \leq d_0(h,\, R(\Pi(h))) = \Delta_R(h).
\end{equation*}
Taking inf over $R$ and sup over $h$:
\begin{equation*}
  \delta(\Pi_2) \leq \delta(\Pi),
\end{equation*}
meaning the defect of the intermediate space $M_1$ is no greater than
that of the full chain. Equivalently, $\delta(\Pi) \geq \delta(\Pi_1)$.
\end{proof}

\begin{corollary}[Monotonicity of Defect]
In a projection chain $M_0 \to M_1 \to \cdots \to O$, the minimum
reconstruction defect is non-decreasing: each further projection can
only increase or maintain the defect. Information is monotonically
lost along any chain.
\end{corollary}

This is the reconstruction-defect version of the data processing
inequality. It establishes irreversibility as monotone along projection
chains: no intermediate step can recover information lost at an earlier
step.

\subsection{The Distinction Functional Along a Chain}

\begin{proposition}[Distinction Decreases Along Chains]
For a projection chain $\Pi = \Pi_2 \circ \Pi_1$, the distinction
between any two histories can only decrease:
\begin{equation}
  \mathcal{D}(\delta_{\Pi(h)}, \delta_{\Pi(h')})
  \leq \mathcal{D}(\delta_{\Pi_1(h)}, \delta_{\Pi_1(h')})
  \leq \mathcal{D}(\delta_h, \delta_{h'}),
\end{equation}
where point masses $\delta_h, \delta_{h'}$ are the degenerate distributions
on $M_0$.
\end{proposition}

\begin{proof}
This follows from the data processing inequality for JSD
(Corollary~\ref{cor:data_proc}): each further projection $\Pi_i$
is a measurable map applied to the distributions, and data processing
ensures $\mathcal{D}(f_*P, f_*Q) \leq \mathcal{D}(P,Q)$.
\end{proof}

The proposition formalises the intuition that passing information through
more processing steps can only blur the distinction between its sources.
A signal that is clearly distinguishable at the physical level may become
indistinguishable after enough stages of perceptual, cognitive, or
institutional processing.

\subsection{When Intermediate Projections Help}

Theorem~\ref{thm:chain_defect} establishes that intermediate projections
cannot decrease defect. But they can \emph{re-structure} it: an
intermediate projection $\Pi_1$ may create a representation $M_1$ from
which subsequent reconstruction $R_1 : O \to M_1$ is more efficient
than direct reconstruction $R : O \to M_0$, even though the total
defect cannot decrease.

This is the formal account of \emph{representation learning}: the task
of finding an intermediate projection $\Pi_1$ such that the downstream
tasks (encoded in $\Pi_2$) can be performed with lower defect from
$M_1$ than from $M_0$ directly. Good representations are intermediate
projections that preserve the distinctions needed for downstream tasks
while collapsing distinctions that are irrelevant to them.

\begin{definition}[Task-Relevant Distinction Preservation]
Given a downstream projection $\Pi_2 : M_1 \to O$, an intermediate
projection $\Pi_1 : M_0 \to M_1$ is \emph{task-relevant} if
$\delta(\Pi_2 \circ \Pi_1) = \delta(\Pi_2)$: the composition has the
same defect as the downstream projection alone. In other words,
$\Pi_1$ preserves all distinctions that matter for $\Pi_2$.
\end{definition}

Task-relevant projections are the ideal representations: they discard
exactly the information that is irrelevant to the downstream task and
preserve exactly the information that is relevant. The theory of
task-relevant projections is the foundation of the theory of abstraction,
compression, and feature learning.


%======================================================================
\section{Conditional Distinctions and Contextual Comparison}
%======================================================================

\subsection{Comparing Interfaces in Context}

The distinction functional $\mathcal{D}(P,Q)$ compares two interfaces
globally: it asks how different the full distributions $P$ and $Q$ are.
But in many applications, what matters is not the global difference
between two interfaces but their difference \emph{in a particular
context}: given that an observation falls in category $C$, how different
are the two interfaces?

\begin{definition}[Conditional Interface]
Given observable distributions $P$ and $Q$ on $O$ and a measurable
subset $C \subseteq O$ with $P(C) > 0$ and $Q(C) > 0$, the
\emph{conditional interfaces} are
\begin{equation}
  P|_C(x) = \frac{P(x)\,\mathbf{1}_{x \in C}}{P(C)},
  \qquad
  Q|_C(x) = \frac{Q(x)\,\mathbf{1}_{x \in C}}{Q(C)}.
\end{equation}
The \emph{conditional distinction functional} is
\begin{equation}
  \mathcal{D}_C(P,Q) = \mathcal{D}(P|_C,\, Q|_C)
  = \operatorname{JSD}(P|_C \| Q|_C).
\end{equation}
\end{definition}

\begin{proposition}[Properties of Conditional Distinction]
\begin{enumerate}[label=\textup{(\roman*)},leftmargin=*]
  \item $\mathcal{D}_C(P,Q) = 0$ iff $P|_C = Q|_C$:
        the interfaces coincide within $C$.
  \item $\mathcal{D}_C(P,Q) \leq \log 2$, with equality iff
        $P|_C$ and $Q|_C$ have disjoint support.
  \item $\mathcal{D}_C(P,Q)$ may exceed or fall below $\mathcal{D}(P,Q)$;
        conditioning can reveal or suppress distinctions.
\end{enumerate}
\end{proposition}

\begin{proof}
Parts (i) and (ii) follow immediately from the corresponding properties
of JSD: $\operatorname{JSD}(P\|Q) = 0 \iff P = Q$, and
$\operatorname{JSD}(P\|Q) \leq \log 2$ with equality iff $P$ and $Q$
have disjoint support. Part (iii) follows by example: if $P$ and $Q$
agree on $C$ but differ outside $C$, then $\mathcal{D}_C(P,Q) = 0 <
\mathcal{D}(P,Q)$; if $P$ and $Q$ agree outside $C$ but differ maximally
within $C$, then $\mathcal{D}_C(P,Q) > \mathcal{D}(P,Q)$.
\end{proof}

\subsection{The Distinction Profile}

Rather than a single global comparison, the \emph{distinction profile}
of two interfaces records how their distinction varies across contexts.

\begin{definition}[Distinction Profile]
Given a partition $\{C_1, \ldots, C_k\}$ of $O$ and distributions $P, Q$
on $O$, the \emph{distinction profile} is the vector
\begin{equation}
  \mathcal{D}_{[C]}(P,Q)
  = \bigl(\mathcal{D}_{C_1}(P,Q),\, \mathcal{D}_{C_2}(P,Q),\,
    \ldots,\, \mathcal{D}_{C_k}(P,Q)\bigr) \in [0, \log 2]^k.
\end{equation}
\end{definition}

The distinction profile records where two theories agree and where they
disagree, at the resolution of the partition $\{C_i\}$. Two theories
with the same global distinction $\mathcal{D}(P,Q)$ may have very
different distinction profiles: one pair may disagree uniformly across
all contexts, while another may agree in most contexts but disagree
sharply in a few.

The distinction profile is the basis for a more refined theory of
falsifiability: rather than asking ``is this theory falsifiable?'' (a
global question), one asks ``in which contexts does this prediction
differ from its alternatives?'' (a context-specific question). Strong
falsifiability is a sharply peaked distinction profile: the two theories
differ significantly in a specific, experimentally accessible context.

\subsection{Mutual Distinction and Shared Information}

\begin{definition}[Mutual Distinction]
Given three distributions $P$, $Q$, and $R$ on $O$, the
\emph{mutual distinction} between $P$ and $Q$ given $R$ is
\begin{equation}
  \mathcal{D}(P, Q \mid R)
  = \mathbb{E}_{x \sim R}\!\left[
      \mathcal{D}(\delta_{P(x)},\, \delta_{Q(x)})
    \right],
\end{equation}
where $\delta_{p}$ is the Bernoulli distribution with parameter $p$.
This measures how much the distinction between $P$ and $Q$ is visible
to a reference observer using distribution $R$.
\end{definition}

When $R = M_{PQ} = \frac{1}{2}(P+Q)$, the mutual distinction reduces
to the distinction functional $\mathcal{D}(P,Q)$. When $R$ is the
distribution of an external observer with different priors or different
resolution, the mutual distinction measures the distinction between $P$
and $Q$ as seen by that observer. This formulation is useful for modeling
partial observability and for the theory of coordinated inquiry: two
agents with different reference distributions $R_1$ and $R_2$ may
perceive the distinction between two theories very differently, even
when observing the same outcomes.


%======================================================================
\section{Application Corollaries}
%======================================================================

The three sections above (§2 on reconstruction defect, §3 on the
distinction functional, and §4 on interface geometry) are the mathematical
core. The following corollaries summarise their application to the
three main domains of the monograph: physics, cognition, and institutions.
Detailed treatments appear in the corresponding parts of \emph{Foundations
of Distinction Geometry}.

\subsection{Physics}

\begin{corollary}[Ghost States]
In a Krein-space quantum field theory, a ghost state $g \in \mathcal{K}_-$
is a hidden admissibility direction (Definition~\ref{def:hidden}) if and
only if $\mathcal{D}(\delta_{\Pi(h)}, \delta_{\Pi(h + \epsilon g)}) = 0$ for all
sufficiently small $\epsilon$. Whether this condition holds in the spin-2
sector of quadratic gravity is the central open problem of the programme.
\end{corollary}

\begin{corollary}[Thermodynamic Arrow]
The thermodynamic arrow of time is the direction in which the minimum
reconstruction defect $\delta(\Pi \circ \phi_t)$ grows under the
coarse-graining projection $\Pi$ from microstates to macrostates, where
$\phi_t$ is the microstate dynamics. The second law is the statement that
this growth is non-negative for measure-preserving $\phi_t$.
\end{corollary}

\begin{corollary}[Equivalent Theories]
Two physical theories are empirically indistinguishable at energy scale
$E$ if and only if $d_{\mathrm{int}}(P^{(E)}_1, P^{(E)}_2) = 0$, where
$P^{(E)}_i$ is the distribution of scattering outcomes at energy $E$
predicted by theory $i$. The Derived Interface Conjecture asserts that
Hilbert-space quantum gravity and Krein-space quadratic gravity satisfy
this condition for all $E \ll M_{\mathrm{Pl}}$.
\end{corollary}

\subsection{Cognition}

\begin{corollary}[Memory Loss]
Cognitive memory loss is the growth of reconstruction defect
$\Delta_R(h_{\mathrm{past}})$ over time, where $h_{\mathrm{past}}$ is
the generating history and $R$ is the reconstruction map from current
memory states to past states. Perfect recall requires injective memory
encoding; all real memory systems have positive reconstruction defect.
\end{corollary}

\begin{corollary}[Learning as Defect Reduction]
Learning is the reduction of task-relevant reconstruction defect: the
update of the projection $\Pi$ such that $\delta(\Pi_{\mathrm{new}}) <
\delta(\Pi_{\mathrm{old}})$ for the task-relevant component of the
defect. A learning event that reduces defect for one task while
increasing it for another is a specialisation; one that reduces defect
across all tasks is a genuine improvement.
\end{corollary}

\begin{corollary}[Paradigm Shift as Projection Change]
A scientific or cognitive paradigm shift is a change of projection
$\Pi \to \Pi'$ that is a distinction reorganisation: $\Pi'$ is neither
a refinement nor a coarsening of $\Pi$. The incommensurability of two
paradigms is measured by $\mathcal{D}(\Pi_*\mu, \Pi'_*\mu)$: how
different the observable interfaces they produce from the same underlying
reality are.
\end{corollary}

\subsection{Institutions}

\begin{corollary}[Goodhart's Law]
Goodhart's Law --- when a measure becomes a target it ceases to be a
good measure --- is the theorem that optimising a proxy $M$ for a true
objective $T$ navigates the fiber $\Pi^{-1}(\text{high } M)$ in
directions that may decrease $T$. The fiber is non-trivial whenever
$\Pi$ is non-injective, which it always is for any realistic measurement
system. The severity of the Goodhart effect is measured by the
distinction $\mathcal{D}(P_M, P_T)$ between the interface of the measure
and the interface of the true objective.
\end{corollary}

\begin{corollary}[Institutional Opacity]
An institution is opaque to degree $\kappa > 0$ at a given parameter if
the sectional curvature of theory space in the direction of the
institution's internal structure is $+\kappa$: many distinct internal
configurations produce nearly identical observable behaviour. High
positive curvature is high institutional opacity.
\end{corollary}

\begin{corollary}[Legibility as Injectivity]
A policy or governance intervention is fully legible if the projection
from the target community's internal dynamics to the policy's observable
outcomes is injective. Scott's catastrophes arise when this projection
is treated as injective (the community is treated as fully legible)
when it is in fact highly non-injective (the community's internal
dynamics are largely invisible to the policy measure).
\end{corollary}


%======================================================================
\section*{Conclusion}
\addcontentsline{toc}{section}{Conclusion}
%======================================================================

This paper has developed three interlocking mathematical structures and
shown that they form a coherent account of what any observation does and
what any comparison of observations reveals.

The reconstruction defect $\Delta_R(h) = d(h, R(\Pi(h)))$ measures the
irreducible cost of non-injective projection: the minimum error any
reconstruction must make, regardless of how cleverly it proceeds. The
Irreversibility from Non-Injectivity theorem establishes that this cost
is always positive when the projection is many-to-one, and the Entropy
Bound connects the minimum defect to the Shannon entropy of the projected
distribution. Irreversibility, memory loss, thermodynamic entropy, and
the arrow of time are all instances of this same structure.

The distinction functional $\mathcal{D}(P,Q) = \operatorname{JSD}(P\|Q)$
measures the information destroyed when the distinction between two
observable interfaces is forgotten. It is interface-native --- defined
without reference to any internal space --- symmetric, and locally
equivalent to the Fisher information metric. Gauge orbits are its
zero-level sets. Hidden admissibility directions are deformations along
which it remains zero. The Born rule of quantum mechanics is the
simplest possible distinction functional: the map that forgets the phase
of a complex amplitude.

The theory space $(\mathcal{M}, d_{\mathrm{int}})$ is the metric space
of observable interfaces, equipped with the Jensen--Shannon distance.
It is locally Riemannian with Fisher metric. Its curvature is the
curvature of distinguishability: positive curvature means overfitting
and opacity; negative curvature means high falsifiability. Geodesics
are paths of minimal accumulated distinction. Theories at distance zero
are empirically indistinguishable.

The unified translation of Theorem~\ref{thm:translation} collects all
the major concepts of physical, cognitive, and institutional theory into
a single framework: gauge orbits, ghost states, thermodynamic
irreversibility, the Born rule, memory loss, Goodhart effects, and
institutional opacity are all instances of the same three-level chain,
\begin{equation*}
  M
  \;\xrightarrow{\;\Pi\;}
  O
  \;\xrightarrow{\;\mathcal{D}\;}
  (\mathcal{M},\, d_{\mathrm{int}}),
\end{equation*}
differing only in what $M$, $O$, and $\Pi$ are in each application.

The paper's argument can be compressed to a single substitution. The
question
\begin{equation*}
  \text{``What exists?''}
\end{equation*}
is replaced by
\begin{equation*}
  \text{``What distinctions survive projection, and with what geometry?''}
\end{equation*}
The distinction functional and the theory space metric are the tools for
answering the second question. They are the tools for doing science,
epistemology, and institutional analysis without presupposing access to
the internal space from which observations are generated.

What remains is the investigation of specific instances. The monograph
\emph{Foundations of Distinction Geometry} develops the physics,
cognitive, and institutional applications in detail. The present paper
has established the mathematical core that those applications draw on.


\begin{thebibliography}{99}

\bibitem{endres2003}
D.~M. Endres and J.~E. Schindelin,
``A new metric for probability distributions,''
\textit{IEEE Transactions on Information Theory} \textbf{49}(7),
1858--1860 (2003).

\bibitem{amari2000}
S.-I. Amari and H.~Nagaoka,
\textit{Methods of Information Geometry},
American Mathematical Society, Providence, 2000.

\bibitem{cover2006}
T.~M. Cover and J.~A. Thomas,
\textit{Elements of Information Theory}, 2nd ed.,
Wiley-Interscience, Hoboken, 2006.

\bibitem{ay2017}
N.~Ay, J.~Jost, H.~V\^an L\^e, and L.~Schwachh\"{o}fer,
\textit{Information Geometry},
Springer, Cham, 2017.

\bibitem{csiszar2004}
I.~Csisz\'{a}r and P.~Shields,
``Information theory and statistics: a tutorial,''
\textit{Foundations and Trends in Communications and Information Theory}
\textbf{1}(4), 417--528 (2004).

\bibitem{nielsen2019}
F.~Nielsen,
``On the Jensen--Shannon symmetrization of distances relying on abstract
means,''
\textit{Entropy} \textbf{21}(5), 485 (2019).

\bibitem{lin1991}
J.~Lin,
``Divergence measures based on the Shannon entropy,''
\textit{IEEE Transactions on Information Theory} \textbf{37}(1),
145--151 (1991).

\bibitem{friedan1980}
D.~Friedan,
``Nonlinear models in $2+\epsilon$ dimensions,''
\textit{Physical Review Letters} \textbf{45}, 1057--1060 (1980).

\bibitem{jaynes2003}
E.~T. Jaynes,
\textit{Probability Theory: The Logic of Science},
Cambridge University Press, Cambridge, 2003.

\bibitem{chentsov1982}
N.~N. Chentsov,
\textit{Statistical Decision Rules and Optimal Inference},
American Mathematical Society, Providence, 1982.

\end{thebibliography}

\end{document}

